One annotator produced the study's labels. This appendix reports a check built
from a source no annotator touched, and the citation machinery behind
\S\ref{sec:long}.

\subsection{Judicial characterizations} When a court writes ``\emph{Smith}, 123
F.3d 456, 460 (declining to decide whether the statute applies),'' it asserts on
the record that \emph{Smith} left that question open, while ``(holding that
\dots)'' asserts the opposite. We scanned all \NEligible{} eligible opinions for reporter
citations and applied a frozen two-class characterization dictionary to a window
around each. \NCharLoose{} passages matched, but the loose form is unreliable, since a
marker near a citation is often about something else, as where ``it did not reach
a verdict'' refers to a jury. Requiring the verb to bind directly to the citation,
as a parenthetical or appositive, leaves \NCharTight{} passages, \NCharLZero{} of them
assertions of non-resolution, over 22,471 cited opinions. We release this set.

\subsection{Limits as a validation source} The obvious use fails. Suppose a later
court says \emph{Smith} held $X$ while our annotation says \emph{Smith} left $X$
open. Two explanations fit, an error in our label, or an overstatement of
\emph{Smith} by the later court, which is the phenomenon this paper measures.
The disagreement is equally consistent with both. We report this because the design
is superficially attractive and the flaw becomes visible only once the counts are
computed.

\subsection{Proximity matching} Pairing each characterization with the sentence in
the cited opinion sharing the most content words gives \NJudgeBench{} items (\NJudgeUnres{}
unresolved, \NJudgeDec{} decided), but systems transfer poorly, with the frozen dictionary
reaching \JudgeLexF{} macro-$F_1$ against \LexF{} on the annotated benchmark. The
cause lies in the matcher, which selects where the opinion \emph{discusses} the
issue rather than where it \emph{disposes} of it.
Asking instead whether a dictionary expression sits within 400 characters of the
characterized passage covers \NEscN{} characterizations, of which \NProxyGood{} of \NProxyCheck{}
hand-checked flags were genuine. Both failures are the same failure, and it is why
\S\ref{sec:long} matches propositions.

\subsection{Chains, matching, and inference} Absolute rates in \S\ref{sec:long}
are contrasts rather than prevalence estimates, since the characterization set is
dominated by holding-ascriptions by construction. The origin-side reading is high
precision and low recall, so the comparison runs between origins that
\emph{visibly} mark an issue open and the rest. Intervals resample originating
opinions rather than passages, since several citing passages share an origin. We
report associations rather than causal estimates, because whether a court leaves
an issue open is a choice
correlated with unobservables, including how contested the question was, and those
plausibly affect how later courts describe it.

Locating where an opinion disposes of a named proposition, rather than where it
discusses one, is the binding constraint on a population rate. At the measured
yields, estimating one to within a few points needs on the order of a thousand
verified propositions, which at one in ten means annotating close to ten thousand
citing passages, an annotation budget an order of magnitude larger than a single
annotator supplies.
